\chapter[Typed Interfaces, Contracts and Safety Supervision]
{Typed Interfaces, Contracts and\\Safety Supervision}
\label{chap:interfaces}

\section{A typed command and state interface}

The typed rail interface separates command/state topics; every message defines
its role, required fields and device, keeping requests distinct from reports.

ROS~2 \code{.msg}, \code{.srv} and \code{.action} files
\citep{ros2jazzyinterfaces} define named state, switch/stopper/shuttle
command--state, sensor and visual messages, plus \code{AddShuttle}, independently
of Gazebo implementation.

\section{Command/state separation}

Three roles are distinct: a \emph{command} is a request; \emph{state} is the
controller report; and a \emph{result} records status, supporting readings and,
where applicable, a reason.

A \emph{command effect} is expected post-publication controller/sensor state,
which the executive observes rather than assumes. It differs from a PDDL
symbolic effect, the planner's predicted change. Exact arrival needs
post-command DZI and controller reports; switch/stopper checks apply their
freshness rules. A report is \emph{fresh} only while its timestamp age is within
the configured limit.

\section{Sparse partial-update semantics}

Live switch/stopper commands are sparse named updates. A \emph{checked command
envelope} is the validated JSON-compatible request before typed-message
conversion. Thus \code{switches: \{A3: INTERIOR\}} changes only A3; omission
means preserve state, not apply a default.

After safety checks, the supervisor publishes \code{SwitchCommand} or
\code{StopperCommand}; their \code{NamedState[]} arrays apply only present
entries. Motion uses \code{ShuttleCommand}. The translator's
\code{event\_action} is non-actuating primitive/target metadata;
training-only \code{target\_mask} marks available label values for loss. Neither
is a live command.

\Cref{fig:contract-semantics} summarises feedback confirmation and named-only
change.

\begin{figure}[htbp]
  \centering
  \resizebox{0.98\textwidth}{!}{\begin{tikzpicture}[
  font=\sffamily\small,
  panel/.style={draw=mfjablue,rounded corners=2pt,inner sep=3mm,
    minimum width=158mm},
  flow/.style={draw,rounded corners=2pt,align=center,minimum height=10mm,
    text width=21mm,fill=mfjalight},
  typed/.style={flow,draw=mfjaorange,fill=mfjaorange!8},
  authority/.style={flow,draw=mfjagreen,fill=mfjagreen!8},
  arr/.style={-{Latex[length=1.9mm]},thick,mfjablue},
  feedback/.style={-{Latex[length=1.9mm]},thick,mfjagreen},
  cell/.style={draw,minimum width=13mm,minimum height=7mm,align=center}
]
  \node[panel,minimum height=66mm] (commandpanel) {};
  \node[anchor=north west,font=\sffamily\bfseries\normalsize,mfjablue]
    at ([xshift=2mm,yshift=-2mm]commandpanel.north west)
    {(a) Command and result are separate};

  \node[flow] (caller) at ([xshift=-57mm,yshift=8mm]commandpanel.center)
    {Task executive\\or operator};
  \node[typed,right=5mm of caller] (command)
    {Checked command\\request};
  \node[authority,right=5mm of command] (supervisor)
    {Command\\supervisor};
  \node[typed,right=5mm of supervisor] (typedcommand)
    {Typed ROS~2\\device message};
  \node[authority,right=5mm of typedcommand] (controller)
    {Responsible\\controller};
  \node[flow,below=10mm of controller] (state)
    {Controller\\state report};
  \node[authority,left=5mm of state] (verify)
    {Compare with\\expected effect};
  \node[flow,left=5mm of verify] (result)
    {Recorded\\command result};
  \node[anchor=north east,draw=mfjared,fill=mfjared!7,
    rounded corners=1.2mm,inner xsep=2.2mm,inner ysep=1.1mm,
    font=\sffamily\bfseries\small,text=mfjared]
    at ([xshift=-3mm,yshift=-3mm]commandpanel.north east)
    {Publication \(\neq\) observed success};
  \draw[arr] (caller) -- (command);
  \draw[arr] (command) -- (supervisor);
  \draw[arr] (supervisor) -- (typedcommand);
  \draw[arr] (typedcommand) -- (controller);
  \draw[feedback] (controller) -- (state);
  \draw[feedback] (state) -- (verify);
  \draw[feedback] (verify) -- (result);

  \node[panel,minimum height=84mm,below=7mm of commandpanel] (maskpanel) {};
  \node[anchor=north west,font=\sffamily\bfseries\normalsize,mfjablue]
    at ([xshift=2mm,yshift=-2mm]maskpanel.north west)
    {(b) A sparse named update changes only the addressed device};

  \node[typed,minimum width=41mm,text width=35mm]
    (sparse) at ([xshift=-48mm,yshift=16mm]maskpanel.center)
    {Envelope\\\texttt{switches: \{A3: INTERIOR\}}};
  \node[authority,minimum width=34mm,text width=28mm,right=8mm of sparse]
    (gate) {Supervisor checks\\the named update};
  \node[typed,minimum width=48mm,text width=42mm,right=8mm of gate]
    (roscommand) {\texttt{SwitchCommand}\\
      \texttt{switches=[}\\\texttt{A3: INTERIOR]}};
  \draw[arr] (sparse) -- (gate);
  \draw[arr] (gate) -- (roscommand);

  \matrix (m) [matrix of nodes,nodes={cell},column sep=-\pgflinewidth,
    row sep=-\pgflinewidth] at ([yshift=-15mm]maskpanel.center) {
    |[draw=none,minimum width=28mm,align=right]| {} &
      |[fill=mfjalight]| {A1} & |[fill=mfjalight]| {A2} &
      |[fill=mfjalight]| {A3} & |[fill=mfjalight]| {A4} \\
    |[draw=none,minimum width=28mm,align=right]| {Current} & E & I & E & I \\
    |[draw=none,minimum width=28mm,align=right,font=\sffamily\bfseries\small]| {Named update} &
      |[fill=black!5,text=black!40]| --- & |[fill=black!5,text=black!40]| --- &
      |[fill=mfjaorange!18]| I & |[fill=black!5,text=black!40]| --- \\
    |[draw=none,minimum width=28mm,align=right,font=\sffamily\bfseries\small]| {Result} &
      |[fill=mfjagreen!12]| E & |[fill=mfjagreen!12]| I & |[fill=mfjagreen!24]| I & |[fill=mfjagreen!12]| I \\
  };
  \node[below=3mm of m,font=\sffamily\bfseries\small,mfjagreen]
    {Only A3 is transmitted and changed; omitted names retain their state.};
\end{tikzpicture}
}
  \caption[Supervised sparse commands and controller feedback.]
  {The supervisor checks sparse commands before typed publication; feedback
  establishes results and omitted devices retain state.}
  \label{fig:contract-semantics}
\end{figure}
\FloatBarrier

\section{Versioned contracts}

Operational/task data use versioned Python contracts: named schemas with
validation, provenance and allowed producers. Records carry a schema version;
incompatible versions are rejected. Planning/execution combine typed records
with checked JSON commands and, where required, identifiers, timestamps and
producer/source data. The supervisor converts authorised requests to typed ROS~2
messages. \Cref{tab:contracts} lists the principal contracts.

\begin{table}[htbp]
\centering
\caption{Principal closed-loop contracts and their producers.}
\label{tab:contracts}
\small
\renewcommand{\arraystretch}{1.04}
\begin{tabularx}{\textwidth}{@{}
  >{\raggedright\arraybackslash}p{0.27\textwidth}
  >{\raggedright\arraybackslash}p{0.34\textwidth}
  >{\raggedright\arraybackslash}X@{}}
\toprule
\textbf{Contract} & \textbf{Meaning} & \textbf{Allowed producer}\\
\midrule
\code{ObservedFact} & Value, confidence and
known/unknown/stale/conflicting status & Visual or rule-based provider,
identified by source\\
\tablerowrule
\code{ObservedState} & Visual inputs and fused facts with freshness/status
metadata & State-fusion layer\\
\tablerowrule
\code{SemanticParseEnvelope} & Strict user-text semantic proposal
& Local language model; semantic fields only, excluding PDDL/plan/command\\
\tablerowrule
\code{TaskGoalDraft} & Incomplete proposal after field precedence/fusion
& Explicit parser and field-fusion resolver\\
\tablerowrule
\code{TaskGoal} & Immutable validated constraints (unchanged after publication),
optionally grounded to one current instance & Domain validator; gateway may
ground from current visual facts\\
\tablerowrule
\code{PddlPlanStep} & One returned-plan symbolic action
& PlanSys2 boundary; only the executive may add a required route-normalisation
prerequisite\\
\tablerowrule
\code{TranslatedPlanStep} & Parsed action, sparse command and event-action
record & Plan translator\\
\tablerowrule
Checked command dictionary & Actuation intent with named updates only
& Executive/macro dispatcher; supervisor-authorised before publication\\
\tablerowrule
\makecell[l]{\code{SwitchCommand}/\\\code{StopperCommand}/\\
\code{ShuttleCommand}} & Typed execution messages; devices carry sparse
\code{NamedState[]} & Rule-based supervisor after current-state validation\\
\tablerowrule
\makecell[l]{\code{PostconditionCheck}\\and closed-loop records} & Effect status,
reason and step/task outcome & Executive using supervisor/controller/observation reports\\
\bottomrule
\end{tabularx}
\end{table}

\emph{Producer} means the component allowed to construct a contract;
\emph{source} is stored fact provenance and controls use. Thus simulator pose
may label or veto but never enters visual-model input. Language proposals require
fusion, validation and operator confirmation. Field fusion combines structured
input, explicit text, confirmed context and semantic proposal by fixed
precedence, recording disagreements. The grounding gateway resolves
\emph{any}/\emph{nearest} against current state; route normalisation restores
the exterior route after special clearance; and a macro dispatcher
deterministically expands an approved macro into checked commands. Producer
boundaries appear in \cref{tab:contracts}; the full sequence is in
\cref{fig:neurosymbolic-architecture}.

\section{Observed-state semantics}

Critical facts preserve four statuses rather than coercing them to false:
\emph{known}; \emph{unknown} (unobserved); \emph{stale} (freshness elapsed); and
\emph{conflicting} (authorised sources disagree). \Cref{fig:observed-fact-eligibility}
shows their planner/actuation eligibility.

\begin{figure}[htbp]
  \centering
  \resizebox{0.98\textwidth}{!}{\begin{tikzpicture}[
  font=\sffamily\small,
  header/.style={draw=mfjablue,rounded corners=2pt,fill=mfjalight,
    align=center,text width=150mm,minimum height=11mm,font=\sffamily\bfseries},
  status/.style={draw,rounded corners=2pt,align=center,text width=39mm,
    minimum height=14mm},
  response/.style={draw=black!30,rounded corners=2pt,fill=black!3,
    align=center,text width=51mm,minimum height=14mm},
  consequence/.style={draw,rounded corners=2pt,align=center,text width=45mm,
    minimum height=14mm,font=\sffamily\bfseries},
  ok/.style={draw=mfjagreen,fill=mfjagreen!8},
  warn/.style={draw=mfjaorange,fill=mfjaorange!8},
  bad/.style={draw=mfjared,fill=mfjared!7},
  edge/.style={-{Latex[length=2mm]},thick,mfjablue},
  reject/.style={-{Latex[length=2mm]},thick,mfjared}
]
  \node[header] (head)
    {{\ttfamily ObservedFact} status check --- exactly one branch is taken};

  \node[status,ok,below=5mm of head,xshift=-56mm] (known)
    {{\color{mfjagreen}\bfseries KNOWN}\\fresh, authorised and consistent};
  \node[response,right=7mm of known] (knownaction)
    {Retained and available\\for planner input};
  \node[consequence,ok,right=7mm of knownaction] (accepted)
    {May populate\\executable planner facts};
  \draw[edge,mfjagreen] (known) -- (knownaction);
  \draw[edge,mfjagreen] (knownaction) -- (accepted);

  \node[status,warn,below=4mm of known] (unknown)
    {{\color{mfjaorange}\bfseries UNKNOWN}\\value was not observed};
  \node[response,right=7mm of unknown] (unknownaction)
    {Recorded; the missing value\\remains explicit};
  \node[consequence,bad,right=7mm of unknownaction] (unknownstop)
    {Not passed\\to the planner};
  \draw[reject] (unknown) -- (unknownaction);
  \draw[reject] (unknownaction) -- (unknownstop);

  \node[status,warn,below=4mm of unknown] (stale)
    {{\color{mfjaorange}\bfseries STALE}\\last reading is too old};
  \node[response,right=7mm of stale] (staleaction)
    {Recorded; a newer\\observation is required};
  \node[consequence,bad,right=7mm of staleaction] (stalestop)
    {Not passed\\to the planner};
  \draw[reject] (stale) -- (staleaction);
  \draw[reject] (staleaction) -- (stalestop);

  \node[status,bad,below=4mm of stale] (conflict)
    {{\color{mfjared}\bfseries CONFLICTING}\\authorised sources disagree};
  \node[response,right=7mm of conflict] (conflictaction)
    {Recorded with the\\source disagreement};
  \node[consequence,bad,right=7mm of conflictaction] (conflictstop)
    {Not passed\\to the planner};
  \draw[reject] (conflict) -- (conflictaction);
  \draw[reject] (conflictaction) -- (conflictstop);

  \node[draw=mfjablue!50,fill=mfjalight,rounded corners=2pt,align=center,
    text width=150mm,minimum height=10mm,below=5mm of conflictaction]
    {Eligible facts enter planning first; translation and supervisor checks\\
      then precede command publication.};
\end{tikzpicture}
}
  \caption[Eligibility gates for facts entering planner state.]
  {Status handling before an observed fact enters planner state.}
  \label{fig:observed-fact-eligibility}
\end{figure}
\FloatBarrier

\code{ObservedState} separates:

\begin{description}
  \item[\code{visual\_model\_inputs}] visual-model-origin facts before source
  fusion; despite the legacy field name, these are learned outputs rather than
  the raw images supplied to inference;
  \item[\code{fused\_planner\_state}] fused facts with their status and
  freshness metadata after source reconciliation.
\end{description}

The live task gateway combines a confirmed task with current state, retaining
the latest accepted visual, supervisor and device reports. A provider supplies
one current executive snapshot. Live operation separates oracle and visual
providers; test-only providers remain outside it.

\section{Supervisor responsibilities}

The rule-based \code{room\_315\_vla\_supervisor.py}---not a learned VLA
policy---gives final software authorisation. It validates sparse JSON commands
against rail state, then publishes typed commands consumed by the kinematic
controller.

The supervisor gates include:

\begin{itemize}
  \item latched emergency-stop state and the immediate global \code{stop\_all}
  operation; command structure, identity and rail-side consistency;
  \item fresh/non-conflicting state, occupied targets, reservation-block
  ownership, geometric separation and block headway;
  \item switch change within the \SI{0.35}{m} protected distance;
  stopper and route-clearance state;
  \item post-command occupancy by another shuttle, sensor dropout, timeout,
  unknown localisation, impossible pose or \code{FALLING} mode;
  \item a configured recovery-retry limit (default two) and lexicographic
  reservation-conflict tie-breaking.
\end{itemize}

Diagnostics retain metrics, recent decisions and each authorisation/rejection
reason; rejection leaves the device unchanged. Exceptions or shutdown during
an active goal trigger gateway \code{stop\_all}.

\section{Reservations and deadlock}

A \emph{reservation block} is a mutually exclusive rail-side/public-segment
route resource, temporarily assigned to one shuttle/task; target-slot conflicts
are separate. Reservations prevent two execution primitives (translated
operations) entering it; geometric separation independently prevents overlap.

Occupancy by a non-owner is a crossed-ownership deadlock conflict. The
lexicographically first identity gains priority; others stop, observe and plan
again. This detector covers crossed conflicts, not temporal non-progress or
larger wait-for cycles (each participant awaiting the next one's resource).
Recovery still needs a planned reachable destination, capacity and effect check.

\section{Source of switch feedback}

The planner uses visual location, while switch safety independently checks
continuous controller coordinates. It maps the public controller segment to
the side-specific internal metric graph and measures distance along rail, not
Euclidean space. Only current, finite, graph-mappable state authorises a switch;
fail-closed handling inhibits all else.

\section{Checking command effects}

The runtime defines the expected effects of every translated execution
command. Examples include:

\begin{itemize}
  \item \code{SET\_SWITCHES}/\code{SET\_STOPPERS}: requested values match
  controller status;
  \item commanded slot motion: two consecutive identity-bearing DZI readings,
  controller
  \path{reached_target_slot}, explicit \code{DISABLED} state and required
  confirmation frames agree;
  \item \code{STOP\_NOW}: controller reports a fresh explicit
  \code{DISABLED} state;
  \item observation-only inspection: a newer visual observation is
  available after the request.
\end{itemize}

Results retain command, plan-step and observed-state identifiers; motion effects
also record command-linked observations and require newer reports for arrival
and retries. Switch/stopper matching permits \emph{cached state}: the last
accepted value may predate the command, unlike post-command motion checks.

\section{Manual and emergency paths}
\enlargethispage{2\baselineskip}

Manual device-debug commands use normal structural/supervisor checks. The
exception is emergency: it immediately invokes all-shuttle \code{stop\_all}
without planner state and latches rejection of motion until explicit clearing.
Stoppers can thus be tested without a language model but remain safety-gated.
